\chapter{Conclusão}\label{6-conclusao}

Na seção~\ref{3-espectro}, apresentamos alguns resultados de espectro na literatura que motivaram a etapa inicial da pesquisa e que geraram alguns resultados iniciais publicados no \emph{XI Encontro de Teoria da computação} (ETC). Inicialmente mostramos que na convexidade $\overrightarrow{P_3}$, o espectro do número de intervalo é contínuo. Com esse resultado, estudamos algumas classes de grafos para calcularmos alguns limitantes inferiores e superiores, sendo elas: Caminhos, Ciclos e Árvores.

Dando continuidade, na seção~\ref{4-algoritmos} apresentamos resultados algorítmicos para calcular o \emph{conjunto de intervalo} e o \emph{conjunto de envoltória} em árvores orientadas. Esse resultado não é uma novidade e sim continuação direta de~\cite{Medeiros2024Convexidade}, onde foi apresentado que existe um algoritmo polinomial para calcular tais parâmetros. Neste caso, apresentamos algoritmos polinomiais explicitos para os respectivos problemas.

Por fim, na seção~\ref{5-complexidade}, apresentamos que calcular o \emph{número de intervalo} $(\overrightarrow{in}_{P_3})$ e o \emph{número de envoltória} $(\overrightarrow{hn}_{P_3})$ em grafos planares, são ambos problemas~\npcs. Em ambos os casos, assim como em~\cite{PENSO201583}, fazemos uma redução do problema~\sam.

O problema do espectro é bem amplo, visto que é possível variar a convexidade e o parâmetro do mesmo. Uma possível linha de trabalhos futuros é exatamente essa, investigar o problema do espectro em outras convexidades, como a convexidade \emph{monofônica}~\cite{duchet1988convex,dourado2010complexity} e possivelmente continuar na convexidade $P_3$ orientada mas para o parâmetro do número de envoltória.

Neste trabalho, foram calculados limitantes para o espectro em classes de grafos bastante conhecidas (ciclos, árvores e caminhos), uma possibilidade é explorar outras classes de grafos como \emph{Splits} ou até diretamente os grafos cordais. Uma outra possibilidade é investigar se os parâmetros do número de intervalo e número de envoltória são tão distantes quanto queira.

Seguindo a linha de complexidade computacional, uma outra possível linha para trabalhos futuros é investigar o número de intervalo em torneios na convexidade $P_3^*$ orientada, ou seja, a convexidade de caminhos orientados de tamanho 2 induzidos. 